\chapter{Related Work}
\label{ch:related-work}

This chapter reviews prior work for the three studies in this thesis. We group it
into three sections. Section~\ref{sec:rw-neurons} covers language-specific neurons
and how earlier work has tried to use internal components to control multilingual
behaviour. Section~\ref{sec:rw-crosslingual} covers how multilingual models
process language across layers, and how this has been studied during and after
RL fine-tuning. Section~\ref{sec:rw-fairpo} covers multi-label
classification, preference optimization, and group-robust optimization, which are
the parts that FairPO builds on.

\section{Language-Specific Neurons and Internal Control of Multilingual Models}
\label{sec:rw-neurons}

\subsection{Identifying Language-Specific Neurons}
\label{sec:rw-neurons-id}

A line of work argues that a small number of neurons in a multilingual model are
responsible for handling a particular language. \citet{tang2024language} propose
Language Activation Probability Entropy (LAPE), which marks a neuron as
language-specific when it activates for one or a few languages and stays quiet for
the rest, and they show that these neurons sit mostly in the first and last layers
of the model. \citet{kojima2024multilingual} study decoder-based models and find
neurons that fire for a single language, with little overlap across languages.
Earlier work on encoder models reached a related conclusion:
\citet{durrani2020analyzing} show that small sets of neurons capture specific
linguistic properties, and \citet{bhattacharya2023unveiling} study language
specificity in the feed-forward layers of transformer models. These methods differ
in how they score and threshold neurons, but they share the assumption that
language is, at least in part, stored in identifiable units.

\subsection{Controlling and Editing Through Neurons}
\label{sec:rw-neurons-control}

Several papers go further and use these neurons to change model behaviour.
\citet{tang2024language} and \citet{kojima2024multilingual} show that turning
language-specific neurons on or off can steer the language a model generates in.
\citet{zhao2024multilingualism} detect language-specific neurons with a method
that needs no labelled data, and report that fine-tuning a small set of them
improves performance in a target language. Related ideas appear in machine
translation and style transfer: \citet{zhu2024landermt} route language-aware
neurons to reduce catastrophic forgetting during translation fine-tuning,
\citet{xie2021importance} allocate neurons between shared and language-specific
knowledge to improve translation, and \citet{lai2024style} steer text style by
turning off source-style neurons. Beyond language, \citet{huo2024mmneuron} study
domain-specific neurons in multimodal models and report accuracy gains from
manipulating them. \citet{zhang2024multilingual} edit factual knowledge across
languages using language-agnostic factual neurons.

\subsection{Limits of Neuron-Level Control}
\label{sec:rw-neurons-limits}

This control is not always reliable. A reason often given is that single neurons
do not correspond to single functions. \citet{elhage2022superposition} describe
superposition, where a neuron takes part in more than one feature, so editing it
affects several behaviours at once. Work that traces facts across languages also
suggests that knowledge is partly shared and partly transferred rather than stored
in separate per-language units \citep{zhao2024tracing}. Most of the steering
results above are reported on generation, that is, on which language the model
writes in, and not on whether the model answers a downstream task correctly. The
study in Chapter~\ref{ch:neurons} \citep{mondal2025language} works in this gap. It
tests whether neuron interventions help on XNLI \citep{conneau2018xnli} and XQuAD
\citep{artetxe2020cross} for low-resource languages, and finds that they do not,
that setting an activation to zero does not turn the neuron off, and that neurons
found for different languages are not independent.

\section{Cross-Lingual Processing in Language Models}
\label{sec:rw-crosslingual}

\subsection{The English-Centric View of Multilingual Processing}
\label{sec:rw-crosslingual-english}

A second group of papers studies how a multilingual model moves information across
its layers when given a non-English input. \citet{wendler2024llamas} track the
hidden states of Llama-2 on a word-translation task and describe three stages: the
hidden states start far from any output token, then pass through a region where
the model can already decode the answer but assigns more probability to its
English form, and finally move to a region specific to the input language. They
read this as evidence that the model works in a space closer to English at middle
layers. \citet{zhao2024multilingualism} propose a similar workflow, in which the
model understands the input, works in English in the middle layers, and produces
output in the target language at the end. \citet{lu2025paths} look at factual
recall and describe a pipeline in which the model uses an English-centric recall
mechanism and then translates the English answer back into the target language;
they point to two failure points, weak use of the English mechanism and wrong
translation back to the target, and propose vector interventions that improve
recall accuracy. Together these papers support the idea that English plays a
special role inside the model, but they study translation and factual recall in
fixed models rather than reasoning during training.

\subsection{Tools for Studying Internal Representations}
\label{sec:rw-crosslingual-tools}

The studies above rely on a few standard tools. The logit lens
\citep{nostalgebraist2020logitlens} reads out the model's prediction at each layer
by applying the output head to intermediate hidden states, which gives a
per-layer view of what the model would predict. To compare representations,
\citet{kornblith2019similarity} introduce Centered Kernel Alignment (CKA), a
similarity measure between two sets of representations that is invariant to
orthogonal transformations and isotropic scaling. CKA has been used to compare
representations across languages in multilingual models
\citep{conneau2020emerging}. The study in Chapter~\ref{ch:grpo-crosslingual} uses
both tools: the logit lens to check whether the model predicts in English at
middle layers, and CKA both across languages and across training checkpoints to
measure how much the representations move during fine-tuning.

\subsection{RL for Reasoning and Its Effect on Languages}
\label{sec:rw-crosslingual-rl}

Reinforcement learning is now a common way to improve reasoning in language
models. Group Relative Policy Optimization (GRPO) \citep{shao2024deepseekmath}
removes the separate value model used in PPO and instead estimates the baseline
from a group of sampled outputs for the same prompt. \citet{guo2025deepseekr1}
use GRPO to train a reasoning model with rewards based only on the correctness of
the final answer. Most of this work is evaluated on English or on a mixed set of
languages without separating the effect per language. Multilingual reasoning
itself is measured on benchmarks such as MGSM \citep{shi2022mgsm}, which
translates grade-school math problems into several languages. A separate line of
work shows that the cause of cross-lingual gaps can lie in the tokeniser:
\citet{petrov2023tokenizer} show that the same text can be split into very
different numbers of tokens across languages, which puts low-resource languages at
a disadvantage before the model runs. Chapter~\ref{ch:grpo-crosslingual} connects
these threads. It fine-tunes a multilingual model on MGSM with GRPO, checks the
English-centric processing claim during training, and finds that whether GRPO can
teach a language tracks tokeniser coverage rather than language identity.

\section{MLC and Preference Optimization}
\label{sec:rw-fairpo}

\subsection{Multi-Label Classification and Class Imbalance}
\label{sec:rw-fairpo-mlc}

Multi-label classification usually trains a model by minimizing an aggregate loss,
most often binary cross-entropy, over all labels at once. This treats every label
the same, which is a problem when some labels are rare or more important than
others. \citet{charte2015imbalance} study class imbalance in multi-label data and
the resampling methods used to address it. \citet{lin2018focal} propose Focal
Loss, which lowers the weight of easy examples so that training focuses on harder
ones, but it still scores each label on its own and does not compare a true label
against a specific competing one. Work on long-tailed multi-label recognition
\citep{guo2021longtailed} and on fairness in multi-label settings
\citep{liu2023simfair} addresses related problems, but does not target a chosen
subset of labels while protecting the rest. FairPO \citep{mondal2025fairpo}, the
study in Chapter~\ref{ch:fairpo}, fills this gap by splitting the labels into a
privileged group to improve and a non-privileged group to maintain. The standard
benchmarks here are MS-COCO \citep{lin2015coco} and NUS-WIDE
\citep{chua2009nuswide}, with a Vision Transformer backbone
\citep{dosovitskiy2021vit}, which FairPO also uses.

\subsection{Preference Optimization}
\label{sec:rw-fairpo-po}

Preference optimization was developed to align language models from preference
pairs. Direct Preference Optimization (DPO) \citep{rafailov2023direct} turns the
reinforcement-learning-from-human-feedback objective into a simple classification
loss over preferred and dispreferred outputs, without a separate reward model.
Several reference-free variants followed. Contrastive Preference Optimization
(CPO) \citep{xu2024contrastive} removes the reference model and adds a
negative-log-likelihood term to keep absolute quality. Simple Preference
Optimization (SimPO) \citep{meng2024simpo} uses the length-normalized average
log-probability as the implicit reward and adds a target margin. These methods
were built for text generation, where a preference is between two full sequences.
FairPO adapts their core idea, preferring a correct option over a confusing one,
to per-label scores in classification, so that the model is trained to score a
true label above the incorrect label most likely to be confused with it.

\subsection{Group-Robust Optimization}
\label{sec:rw-fairpo-dro}

The last component is the method used to balance the two groups. Distributionally
Robust Optimization, and in particular Group DRO \citep{sagawa2020dro}, minimizes
the worst-case loss over a set of predefined groups instead of the average loss,
which improves performance on the weakest group. Group Robust Preference
Optimization \citep{ramesh2024group} brings this idea to preference learning,
balancing performance across preference groups with adaptive weights in a minimax
setup. FairPO uses this formulation, but defines the groups by the label partition
rather than by preference source, and balances the preference-based loss for the
privileged group against the constrained loss for the non-privileged group. We
note again that this Group Robust Preference Optimization is a different method
from the Group Relative Policy Optimization of
Section~\ref{sec:rw-crosslingual-rl}, despite the shared abbreviation.